% ============================================================================
% 04_framework.tex — Empirical framework
% ============================================================================

\section{Empirical framework}
\label{sec:framework}

The simulation compares two housing tenure strategies on a cash-flow-matched basis: both start with identical capital, and their cash flows are equalised every month. Whenever the renter pays less for housing in a given month, the difference is invested; whenever the renter pays more, the difference is withdrawn from the portfolio. The design follows \textcite{beracha2012}, with transaction costs and taxes calibrated to the English first-time buyer.

\subsection{The buyer}

The buyer purchases the property at the market price prevailing in the entry month. In the baseline scenario, they put down a deposit of 10\% of the price (loan-to-value of 90\%), pay purchase costs of 1.2\% of the price (conveyancing, survey, mortgage and valuation fees), and stamp duty under the first-time-buyer schedule in force in that month. The mortgage amortises over 30 years.

The baseline rate is a five-year fix, because 99\% of first-time-buyer lending is fixed-rate and the five-year band is the modal choice. The payment is set at entry from the quoted five-year rate at the buyer's loan-to-value, held constant for 60 months, then reset at the rate prevailing at the refix:
\begin{equation}
  r_{m,t} = (1 + r_{t})^{1/12} - 1, \qquad
  \text{Payment}_t = \frac{B_t\, r_{m,t}}{1 - (1 + r_{m,t})^{-(T-t)}},
  \label{eq:payment}
\end{equation}
where $r_t$ is the annual rate in force, $B_t$ the outstanding balance and $T-t$ the remaining months. At each refix the rate is taken from the band matching the loan-to-value the buyer has amortised down to, so a borrower who started at 90\% is not charged the high-loan-to-value premium indefinitely. Appendix~\ref{app:data} sets out the rate construction. A floating rate and a two-year fix are examined in Section~\ref{sec:robustness}, where the fixation design turns out to move the answer more than any parameter except maintenance.

Each month the owner also pays maintenance and depreciation of 1.5\% per year of the property's current market value,
\begin{equation}
  M_t = \frac{0.015}{12}\, V_t,
  \label{eq:maintenance}
\end{equation}
an all-in charge covering routine maintenance, major works and replacements, buildings insurance, and physical depreciation not offset by spending. It is the single most consequential assumption in the paper, so it is built from the national accounts rather than from convention.

The charge has two components. \emph{Depreciation} is measured by ONS consumption of fixed capital on dwellings, divided by the market value of the dwelling stock including the land beneath it: 1.20\% per year on average over 1995--2024, 1.28\% in 2024, with a range of 0.95\% to 1.47\% \parencite{ons_cfc,ons_balance_sheet}. \emph{Cash maintenance and buildings insurance} add roughly 0.35\%, from ONS Family Spending \parencite{ons_family_spending}. Together, 1.55\%, which the baseline rounds to 1.5\%.

The land share is what makes this figure specific to the United Kingdom, and it is why figures imported from the American literature do not transfer. Depreciation applies to structures, not to the ground they stand on, and in the UK land is 69.9\% of the market value of household residential property. A rate estimated on structure value, or on total value in a market where land is half of it, therefore overstates the charge here by a wide margin once applied to a British price. On a UK basis, ONS implies depreciation of 4.24\% of \emph{structure} value, which is 1.28\% of market value at a 30.1\% structure share.

Section~\ref{sec:robustness} reports 1.2\% to 1.8\% as the range UK evidence supports, the floor being depreciation alone, for a buyer who does no discretionary upkeep, and the ceiling being the peak observed depreciation year plus cash maintenance. Values of 1\% and 2--3\% are also reported, but as bounds beyond the evidence rather than as equally supported alternatives. The property value itself evolves with the relevant house-price index, $V_{t+1} = V_t (1 + g_t)$.

The buyer's terminal wealth after a holding period ending at $T^{*}$ is net housing equity after selling costs of 1.8\% (estate-agent commission, conveyancing, energy certificate):
\begin{equation}
  W^{\text{own}} = V_{T^{*}} (1 - 0.018) - B_{T^{*}}.
  \label{eq:nw_buy}
\end{equation}

\subsection{The renter}

The renter occupies an equivalent dwelling at the market rent prevailing each month and invests all spare cash. The initial portfolio equals the buyer's total initial outlay (deposit plus purchase costs plus stamp duty), and each month the difference between the owner's outflow (mortgage payment plus maintenance) and the rent is added to, or withdrawn from, the portfolio, which earns the monthly return on the global equity index (net of withholding taxes and fund fees). The renter's terminal wealth is the portfolio value, $W^{\text{rent}} = P_{T^{*}}$. 

\subsection{Outcome measure and analyses}

Outcomes are summarised by the renter/owner wealth ratio
\begin{equation}
  R = \frac{W^{\text{rent}}}{W^{\text{own}}},
  \label{eq:ratio}
\end{equation}
following \textcite{felix2025}; $R > 1$ means renting-and-investing produced greater terminal wealth. The ratio is invariant to the initial capital scale and comparable across cohorts and regions.

Four analyses are run: (i) 199 entry cohorts, one for every month from January 2005 to July 2021, each holding for exactly five years, which is the primary result; (ii) the same exercise repeated at every holding period from one to ten years; (iii) both region by region, aggregated with population weights; and (iv) the forward-looking required appreciation rate of Section~\ref{sec:rar}, at end-of-sample conditions.

The horizon panel stops at ten years because the sample cannot support more. A fifteen-year hold can only be evaluated for cohorts entering in the first seventy-nine months and a twenty-year hold only for the first nineteen, all of them buying into the pre-crisis peak, so those figures would describe one entry window rather than the effect of holding longer. Ten years still rests on 139 cohorts spanning eleven and a half years of entry dates.

Cohorts start every month rather than every quarter, which removes an arbitrary grid at no cost. It should not be mistaken for extra evidence: the underlying sample is the same 258 months either way, and it contains only 4.3 non-overlapping five-year windows however often the simulation is restarted. Cohort counts in this paper are counts of overlapping simulations, not of independent observations, and Section~\ref{sec:predictors} takes that seriously where it matters.

\subsection{Stamp duty}

Stamp Duty Land Tax is modelled period-accurately: each simulated purchase pays the tax implied by the England first-time-buyer schedule in force in that month, including the 2008--2009 holiday, the 2010--2012 first-time-buyer relief, the 2014 slab-to-slice reform, the permanent relief from November 2017, the 2020--2021 holiday, and the 2022--2025 threshold changes \parencite{govuk_sdlt,hoc_sdlt}. Appendix~\ref{app:sdlt} tabulates the schedules. For the representative dwelling, a first-time buyer paid SDLT in 143 of 258 purchase months (55\%), averaging \gbp{1{,}825} (0.88\% of the price) when payable and never exceeding \gbp{3{,}144}. Stamp duty is therefore a second-order factor nationally. It is not second-order in London, where the average dwelling now sits above the \gbp{500{,}000} ceiling at which first-time-buyer relief is withdrawn entirely, and the bill jumps to \gbp{17{,}694} (Section~\ref{sec:rar}).

\subsection{Scope and abstractions}

The model excludes several factors by design. Income and capital-gains taxes are omitted on both sides: a first-time buyer's main residence is exempt from capital-gains tax, and the renter's flows fit within annual ISA allowances for most of the sample, so the symmetric no-tax treatment is a reasonable approximation for the typical household (though not for very large portfolios held outside tax wrappers). Council tax and utilities are paid by the occupant under both tenure options and net out. In the baseline the renter never borrows to invest, so leverage is available only through the mortgage; Section~\ref{sec:robustness} relaxes that and lets the renter borrow on margin. Non-financial considerations, among them security of tenure, freedom to modify, flexibility to move and the insurance value of ownership against local rent risk \parencite{sinai2005}, are outside the accounting and are discussed in Section~\ref{sec:discussion}. Table~\ref{tab:assumptions} collects the baseline parameters.

Two things about that table are easy to mistake for claims about the world.

First, \textbf{every value in it is a baseline and not a prediction about any individual household}. Each has been calibrated to the best available evidence on what an English first-time buyer typically does: the deposit share and term to Financial Conduct Authority loan-level data, the fixation design to the regulator's published breakdown of fixed-rate lending, the transaction costs to first-time-buyer fee surveys, the maintenance charge to national-accounts depreciation. But a typical value is not a universal one. Deposits range from 5\% to 25\%, terms from 25 to 40 years, and a buyer's own figures may sit anywhere in those ranges. Section~\ref{sec:robustness} reports how the result moves as each parameter varies across its defensible range.

Second, the parameters are not equally consequential. The maintenance charge, the fixation design and the renter's portfolio move the answer a great deal; the transaction costs and the mortgage term barely move it. Section~\ref{sec:robustness} ranks them.

\begin{table}[H]
  \centering
  \caption{Baseline parameters (first-time buyer)}
  \label{tab:assumptions}
  \begin{threeparttable}
    \begin{tabularx}{\textwidth}{@{}lXr@{}}
      \toprule
      \textbf{Parameter} & \textbf{Description} & \textbf{Value} \\
      \midrule
      Property              & Representative English dwelling (Section~\ref{sec:representative}) & --- \\
      Holding period        & Typical tenure in a first home                     & 5 years \\
      Deposit               & Share of purchase price (90\% loan-to-value)       & 10\% \\
      Mortgage term         & Amortisation period                                & 30 years \\
      Rate type             & Five-year fix, refixed at the prevailing five-year rate, repriced by current loan-to-value & --- \\
      Purchase costs        & Conveyancing, survey, mortgage product fee         & 1.2\% \\
      Stamp duty            & England first-time-buyer schedule, by month        & rule-based \\
      Selling costs         & Estate agent, conveyancing, EPC                    & 1.8\% \\
      Maintenance \& depreciation & Share of current market value, per year      & 1.5\% \\
      Equity index          & MSCI ACWI net total return, GBP, less 0.12\% fee   & --- \\
      Renter's leverage     & Unlevered portfolio                                & 1.0$\times$ \\
      Rent                  & Market rent, updated monthly                       & --- \\
      Currency              & Nominal GBP throughout                             & --- \\
      \bottomrule
    \end{tabularx}
  \end{threeparttable}
\end{table}
